\ClassNoteNoLine{\ClassName}{Fontset key is \mit@fontset}
\IfFileExists{fontsets/\ClassName-\mit@fontset.tex}{% 2023/07/03
	\input{fontsets/\ClassName-\mit@fontset.tex}
}{%
	\ClassWarning{\ClassName}{Fontset file \ClassName-\mit@fontset.tex or subdirectory fontsets not found, will look for \mit@fontset.tex in your working directory}
	\IfFileExists{\mit@fontset.tex}{% 2023/08/30
		\input{\mit@fontset.tex}
		}{
    	\ClassWarning{\ClassName}{Fontset file \mit@fontset.tex not found, using default fonts. You may need to place the fontset file into your working directory}
    	\ifpdftex
        	\ClassNoteNoLine{\ClassName}{Loading Computer Modern Roman text and math fonts (default for pdftex)}
        	\RequirePackage[T1]{fontenc}
	        \RequirePackage{amssymb}% covered by unicode-math, libertinus, lucidabr, newtx, newtxsf packages.
        	\RequirePackage{bm}	
        	% Patch headings to automatically use bold math with CMR fonts. Compare to report.cls. 
            % Predefined fontsets for which bold math exists: defaultfonts, fira-newtxsf, newtx, newtx-sans-text (all 4 in pdf only); + Lucida.
            % Those fontset files also include this code. 
        	\patchcmd{\section}{\normalfont}{\mathversion{bold}\normalfont}{}{}
        	\patchcmd{\subsection}{\normalfont}{\mathversion{bold}\normalfont}{}{}
        	\patchcmd{\subsubsection}{\normalfont}{\mathversion{bold}\normalfont}{}{}
        	\patchcmd{\paragraph}{\normalfont}{\mathversion{bold}\normalfont}{}{}
        	\patchcmd{\subparagraph}{\normalfont}{\mathversion{bold}\normalfont}{}{}
        \else
        	\ClassNoteNoLine{\ClassName}{Loading Latin Modern text and math fonts (default for unicode engines)}
            \RequirePackage[warnings-off={mathtools-colon,mathtools-overbracket},mathbf=sym]{unicode-math}
            % suppress warnings about lack of integration between mathtools and unicode-math.
            % add mathbf=sym for proper html rendering, 2026/05/12
            % Unicode-math loads fontspec package, but default fonts are not set up by fontspec in the default case.
        \fi
       }
}